\documentclass[aspectratio=169,11pt]{beamer}
% Shared CMPSC 480 Beamer preamble.
\usepackage[T1]{fontenc}
\usepackage[utf8]{inputenc}
\usepackage{lmodern}
\usepackage{amsmath,amssymb,mathtools}
\usepackage{booktabs,array,tabularx}
\usepackage{xcolor}
\usepackage{listings}
\usepackage{tikz}
\usetikzlibrary{arrows.meta,positioning,shapes.geometric}

% Ignore only negligible TeX box diagnostics that are below visual tolerance.
% Larger layout defects still appear in the build log and in rendered QA.
\vfuzz=3pt
\hbadness=2000

\definecolor{courseink}{HTML}{111111}
\definecolor{coursegray}{HTML}{EDEDED}
\definecolor{courserule}{HTML}{B8BCC4}
\definecolor{courseblue}{HTML}{3D8DFF}
\definecolor{courselightblue}{HTML}{D0EDFA}
\definecolor{coursemuted}{HTML}{5C6470}
\definecolor{coursegreen}{HTML}{0B7A53}
\definecolor{coursered}{HTML}{B3261E}

\usetheme{default}
\usefonttheme{professionalfonts}
\setbeamertemplate{navigation symbols}{}
\setbeamersize{text margin left=0.65cm,text margin right=0.65cm}

\setbeamercolor{normal text}{fg=courseink,bg=white}
\setbeamercolor{frametitle}{fg=courseink,bg=white}
\setbeamercolor{title}{fg=courseink,bg=white}
\setbeamercolor{subtitle}{fg=coursemuted,bg=white}
\setbeamercolor{structure}{fg=courseblue}
\setbeamercolor{alerted text}{fg=coursered}
\setbeamercolor{block title}{fg=courseink,bg=courselightblue}
\setbeamercolor{block body}{fg=courseink,bg=coursegray}
\setbeamercolor{example text}{fg=coursegreen}

\setbeamerfont{title}{size=\fontsize{25}{29}\selectfont,series=\bfseries}
\setbeamerfont{subtitle}{size=\fontsize{13}{16}\selectfont}
\setbeamerfont{frametitle}{size=\fontsize{19}{22}\selectfont,series=\bfseries}
\setbeamerfont{normal text}{size=\fontsize{11}{14}\selectfont}
\setbeamerfont{footline}{size=\fontsize{7.5}{9}\selectfont}

\setbeamertemplate{frametitle}{%
  \nointerlineskip
  % Use content-driven height so a long assessment title can wrap without
  % being clipped by a fixed-height title bar.
  \begin{beamercolorbox}[wd=\paperwidth,sep=0.17cm,leftskip=0.48cm,rightskip=0.48cm]{frametitle}
    \insertframetitle
  \end{beamercolorbox}
  \vspace{-0.08cm}
  {\color{courserule}\rule{\textwidth}{0.35pt}}
}

\setbeamertemplate{footline}{%
  \leavevmode
  \begin{beamercolorbox}[wd=\paperwidth,ht=2.6ex,dp=1.1ex,leftskip=.65cm,rightskip=.65cm]{author in head/foot}
    \color{coursemuted}\insertshorttitle\hfill\insertframenumber/\inserttotalframenumber
  \end{beamercolorbox}
}

\setbeamertemplate{itemize item}{\color{courseblue}\small$\blacktriangleright$}
\setbeamertemplate{itemize subitem}{\color{courseblue}\scriptsize$\bullet$}
\setbeamertemplate{enumerate item}{\color{courseblue}\insertenumlabel.}

\lstdefinestyle{coursepython}{
  language=Python,
  basicstyle=\ttfamily\fontsize{8.5}{10}\selectfont,
  keywordstyle=\color{courseblue}\bfseries,
  commentstyle=\color{coursemuted}\itshape,
  stringstyle=\color{coursegreen},
  showstringspaces=false,
  columns=fullflexible,
  keepspaces=true,
  frame=single,
  rulecolor=\color{courserule},
  backgroundcolor=\color{coursegray},
  breaklines=true,
  tabsize=4,
  xleftmargin=0.15cm,
  xrightmargin=0.15cm,
  framexleftmargin=0.15cm,
  framexrightmargin=0.15cm
}
\lstset{style=coursepython}

\newcommand{\points}[1]{\hfill{\color{courseblue}\textbf{[#1 points]}}}
\newcommand{\deliverable}[1]{\textbf{\color{courseblue}#1}}
\newcommand{\code}[1]{\texttt{#1}}
\newcommand{\keyidea}[1]{%
  \begin{beamercolorbox}[rounded=false,sep=0.55em,wd=\linewidth]{block body}
    \textbf{Key idea.} #1
  \end{beamercolorbox}
}
\newcommand{\answerlabel}{\textbf{\color{coursegreen}Answer.}\ }
\newcommand{\gradinglabel}{\textbf{\color{courseblue}Grading guidance.}\ }

\newenvironment{questionblock}[1]
  {\begin{block}{#1}}
  {\end{block}}

\newenvironment{solutionblock}[1]
  {\begin{block}{#1}}
  {\end{block}}

\AtBeginSection[]{
  \begin{frame}
    \vfill
    {\usebeamerfont{title}\color{courseink}\insertsectionhead\par}
    \vspace{0.4cm}
    {\color{courseblue}\rule{0.32\paperwidth}{3pt}}
    \vfill
  \end{frame}
}


% CMPSC480_TOTAL_POINTS: 20
% CMPSC480_ITEM_IDS: 1,2,3,4,5
\title[Module Project B Questions]{Week 7 Module Project B: Q-Learning with Bayesian Belief Tracking}
\subtitle{Acting under partial observability with normalized posterior state}
\author{CMPSC 480 - Student Question Deck}
\date{}

\begin{document}


\begin{frame}
  \titlepage
\end{frame}

\begin{frame}{Assessment overview}
\begin{columns}[T,onlytextwidth]
\begin{column}{0.62\textwidth}
\Large Integrate a from-scratch Bayesian hazard belief with a Q-learning policy.

\vspace{0.35cm}
\normalsize
Coverage: Weeks 4-7: MDPs, Q-learning, approximation, Bayes nets, and inference
\end{column}
\begin{column}{0.33\textwidth}
\begin{block}{Points}20 total\end{block}
\begin{block}{Expected time}10-14 hours\end{block}
\begin{block}{Submission}Repository, evaluation report, and belief traces\end{block}
\end{column}
\end{columns}
\end{frame}

\begin{frame}{Learning objectives}
\begin{itemize}
  \item Update a categorical belief using Bayes rule.
  \item Handle impossible or numerically tiny evidence defensively.
  \item Expose belief summaries to a Q-learning state representation.
  \item Train and evaluate under several sensor-noise regimes.
  \item Ablate belief tracking against a nonbelief baseline.
\end{itemize}
\end{frame}

\begin{frame}{Requirements checklist}
\small
\begin{itemize}
  \item Use Python 3.11 and NumPy; do not use a library that implements the assessed algorithm.
  \item Keep data, model, optimization, and evaluation responsibilities behind explicit interfaces.
  \item Validate shapes, ranges, finite values, empty inputs, and terminal or degenerate cases.
  \item Use deterministic seeds and record every configuration used for reported evidence.
  \item Include focused tests and a README with exact reproduction commands.
\end{itemize}
\end{frame}

\begin{frame}{Task 1 - Belief model (5 points)}
\small
Implement a categorical prior, sensor likelihood, and normalized posterior over candidate hazard locations.

\vspace{0.2cm}
\begin{enumerate}
\item[\textbf{a.}] Write the vectorized posterior update. \hfill{\color{courseblue}\textbf{[2 points]}}
\item[\textbf{b.}] Hand compute prior [0.5,0.5] with likelihoods [0.8,0.2] for one observation. \hfill{\color{courseblue}\textbf{[2 points]}}
\item[\textbf{c.}] Define behavior when the normalization constant is zero or nonfinite. \hfill{\color{courseblue}\textbf{[1 points]}}
\end{enumerate}
\end{frame}

\begin{frame}{Task 2 - RL agent (4 points)}
\small
Implement terminal-aware Q-learning using only information available to the agent.

\vspace{0.2cm}
\begin{enumerate}
\item[\textbf{a.}] Define the observable agent state and prove it excludes the true hidden hazard. \hfill{\color{courseblue}\textbf{[2 points]}}
\item[\textbf{b.}] Specify epsilon-greedy selection and the terminal Q update. \hfill{\color{courseblue}\textbf{[2 points]}}
\end{enumerate}
\end{frame}

\begin{frame}{Task 3 - Integration contract (4 points)}
\small
Connect observations to posterior updates and posterior state to action selection.

\vspace{0.2cm}
\begin{enumerate}
\item[\textbf{a.}] State the within-step order from observation to action to learning. \hfill{\color{courseblue}\textbf{[2 points]}}
\item[\textbf{b.}] Test that reset restores the declared prior and clears episode-local state. \hfill{\color{courseblue}\textbf{[1 points]}}
\item[\textbf{c.}] List every component needed to resume training exactly. \hfill{\color{courseblue}\textbf{[1 points]}}
\end{enumerate}
\end{frame}

\begin{frame}{Task 4 - Noise study and ablation (4 points)}
\small
Evaluate policy safety and reward at low, medium, and high sensor noise.

\vspace{0.2cm}
\begin{enumerate}
\item[\textbf{a.}] Define safety and performance metrics and run multiple seeds. \hfill{\color{courseblue}\textbf{[2 points]}}
\item[\textbf{b.}] Ablate the posterior input and interpret the difference. \hfill{\color{courseblue}\textbf{[2 points]}}
\end{enumerate}
\end{frame}

\begin{frame}{Task 5 - Defensive evidence (3 points)}
\small
Test belief and policy behavior on adversarial boundary cases.

\vspace{0.2cm}
\begin{enumerate}
\item[\textbf{a.}] Provide four tests and name the defect each exposes. \hfill{\color{courseblue}\textbf{[2 points]}}
\item[\textbf{b.}] Explain why a high-noise posterior should remain closer to the prior. \hfill{\color{courseblue}\textbf{[1 points]}}
\end{enumerate}
\end{frame}

\begin{frame}{Edge cases and defensive programming}
\small
\begin{itemize}
  \item Prior has a negative entry.
  \item Evidence has zero likelihood under every hypothesis.
  \item Sensor is completely uninformative.
  \item Episode begins terminal.
  \item True hidden state accidentally enters features.
\end{itemize}
\end{frame}

\begin{frame}{Submission instructions}
\small
\begin{itemize}
  \item Submit source, tests, README, configuration, and the requested report.
  \item Provide one command that reproduces the primary result from a clean environment.
  \item Exclude credentials, generated caches, and unlicensed data.
\end{itemize}
\vspace{0.2cm}
\alert{Do not submit credentials, generated caches, or unapproved libraries.}
\end{frame}

\begin{frame}{Grading rubric - part 1}
\scriptsize
\begin{tabularx}{\textwidth}{>{\bfseries}p{0.28\textwidth} p{0.08\textwidth} X}
\toprule
Category & Points & Required evidence \\
\midrule
Belief model & 5 & BeliefTracker, fixed-sequence oracle tests, and trace table. \\
RL agent & 4 & Agent interface, update tests, and reproducible training loop. \\
Integration contract & 4 & End-to-end sequence diagram, integration tests, and checkpoint serialization. \\
\bottomrule
\end{tabularx}
\end{frame}

\begin{frame}{Grading rubric - part 2}
\scriptsize
\begin{tabularx}{\textwidth}{>{\bfseries}p{0.28\textwidth} p{0.08\textwidth} X}
\toprule
Category & Points & Required evidence \\
\midrule
Noise study and ablation & 4 & Controlled result table with belief and no-belief baselines. \\
Defensive evidence & 3 & Focused tests and a short failure analysis. \\
\bottomrule
\end{tabularx}
\end{frame}

\begin{frame}{Reflection prompt}
In 150-250 words:

Describe one case where the belief was mathematically correct but the policy still behaved poorly. Separate inference error from control or learning error.
\end{frame}

\end{document}
